\begin{theorem}[残余歧义观测的完全谱分解定理]\label{thm:conclusion-statnull-ambiguity-observable-spectral-decomposition}
\leanverified{paper\_conclusion\_statnull\_ambiguity\_observable\_spectral\_decomposition}
固定分辨率 \(m\)。对任意长度 \(n\) 的合法宏观前缀 \(w\in X_m^n\)，记确定化滤波状态的候选数为
\[
K_n(w):=|S_n(w)|.
\]
对任意单调函数
\[
\Phi:\{1,\dots,|V_m|\}\to\RR
\]
定义熵标度归一化的歧义观测量
\[
\mathcal A_{\Phi,m}(n):=
2^{-n}\sum_{w\in X_m^n}\bigl(\Phi(K_n(w))-\Phi(1)\bigr),
\]
并记离散增量
\[
\Delta\Phi(k):=\Phi(k)-\Phi(k-1)\qquad (k\ge 2).
\]
则有精确恒等式
\[
\mathcal A_{\Phi,m}(n)
=
\sum_{k=2}^{|V_m|}\Delta\Phi(k)\,\frac{N_{m,\ge k}(n)}{2^n},
\]
其中
\[
N_{m,\ge k}(n):=\#\{\,w\in X_m^n:\ |S_n(w)|\ge k\,\}.
\]
若对一切满足 \(\Delta\Phi(k)>0\) 的 \(k\)，诱导子图 \(H_{m,\ge k}\) 都满足定理 \ref{thm:Ym-ambiguity-spectrum-law} 的“唯一最大谱半径强连通分量且原始”假设，则
\[
\mathcal A_{\Phi,m}(n)
=
L_{\Phi,m}e^{-\varepsilon_{\Phi,m}n}(1+o(1)),
\]
其中
\[
\varepsilon_{\Phi,m}:=
\min\{\,\varepsilon_{m,\ge k}:\ \Delta\Phi(k)>0\,\},
\]
\[
L_{\Phi,m}:=
\sum_{\substack{2\le k\le |V_m|\\ \Delta\Phi(k)>0\\ \varepsilon_{m,\ge k}=\varepsilon_{\Phi,m}}}
\Delta\Phi(k)\,c_{m,\ge k}>0.
\]
亦即，任意单调清晰度观测在大时间端都完全由阈值歧义谱 \(\{\varepsilon_{m,\ge k}\}_{k\ge 2}\) 支配，其首指数无需重新计算，只需在被激活的阈值层中取最小者。
\end{theorem}
\begin{proof}
对任意整数 \(s\ge 1\)，离散望远镜恒等式给出
\[
\Phi(s)-\Phi(1)
=
\sum_{k=2}^{s}\Delta\Phi(k)
=
\sum_{k=2}^{|V_m|}\Delta\Phi(k)\,\mathbf 1_{\{s\ge k\}}.
\]
代入 \(s=K_n(w)\) 并对全部 \(w\in X_m^n\) 求和，可得
\[
\mathcal A_{\Phi,m}(n)
=
\sum_{k=2}^{|V_m|}\Delta\Phi(k)\,\frac{N_{m,\ge k}(n)}{2^n}.
\]
若对所有 \(\Delta\Phi(k)>0\) 的层级都可应用定理 \ref{thm:Ym-ambiguity-spectrum-law}，则
\[
\frac{N_{m,\ge k}(n)}{2^n}
=
c_{m,\ge k}e^{-\varepsilon_{m,\ge k}n}
+o\!\bigl(e^{-\varepsilon_{m,\ge k}n}\bigr).
\]
由于求和仅有有限项，取其中指数最小者即得首项
\[
L_{\Phi,m}e^{-\varepsilon_{\Phi,m}n},
\]
其余项都以更快指数衰减。证毕。
\end{proof}

\begin{corollary}[精确歧义层的暴露相定理]\label{cor:conclusion-exact-k-ambiguity-exposed-phase}
\leanverified{paper\_conclusion\_exact\_k\_ambiguity\_exposed\_phase}
对每个 \(k\ge 2\)，定义“恰有 \(k\) 个候选”的前缀计数
\[
N_{m,=k}(n):=\#\{\,w\in X_m^n:\ |S_n(w)|=k\,\}.
\]
则恒有
\[
N_{m,=k}(n)=N_{m,\ge k}(n)-N_{m,\ge k+1}(n).
\]
若某一级阈值谱出现严格间隙
\[
\varepsilon_{m,\ge k}<\varepsilon_{m,\ge k+1},
\]
则
\[
\frac{N_{m,=k}(n)}{2^n}
=
c_{m,\ge k}e^{-\varepsilon_{m,\ge k}n}(1+o(1)).
\]
因此，每一个严格的谱台阶都会把“恰好 \(k\) 重歧义”暴露成一个独立指数相，而不再只是尾事件 \(\{|S_n|\ge k\}\) 的粗粒度影子。
\end{corollary}
\begin{proof}
恒等式
\[
N_{m,=k}(n)=N_{m,\ge k}(n)-N_{m,\ge k+1}(n)
\]
显然成立。再由定理 \ref{thm:Ym-ambiguity-spectrum-law}，
\[
\frac{N_{m,\ge k}(n)}{2^n}
=
c_{m,\ge k}e^{-\varepsilon_{m,\ge k}n}(1+o(1)),
\]
\[
\frac{N_{m,\ge k+1}(n)}{2^n}
=
c_{m,\ge k+1}e^{-\varepsilon_{m,\ge k+1}n}(1+o(1)).
\]
若 \(\varepsilon_{m,\ge k}<\varepsilon_{m,\ge k+1}\)，则第二项相对第一项指数级可忽略，故结论成立。证毕。
\end{proof}

\begin{theorem}[阈值歧义集的 Hausdorff 余维阶梯定理]\label{thm:conclusion-threshold-ambiguity-hausdorff-codim-ladder}
对每个 \(k\ge 2\)，定义阈值歧义集
\[
\mathcal Y_{m,\ge k}:=
\{\,y\in Y_m:\ |S_n(y)|\ge k\ \text{对一切 }n\ge 1\,\}.
\]
在前缀超度量
\[
d(y,y'):=2^{-\min\{\,n\ge 0:\ y_n\neq y_n'\,\}}
\]
下，若 \(H_{m,\ge k}\neq\varnothing\) 且满足定理 \ref{thm:Ym-ambiguity-spectrum-law} 的 Perron--Frobenius 假设，则
\[
\dim_{\mathrm H}(\mathcal Y_{m,\ge k})
=
\frac{\log\rho_{m,\ge k}}{\log 2}
=
1-\frac{\varepsilon_{m,\ge k}}{\log 2},
\]
从而
\[
\mathrm{codim}(\mathcal Y_{m,\ge k})
=
\frac{\varepsilon_{m,\ge k}}{\log 2}.
\]
因此，歧义并不只对应一个总余维，而是天然分裂成一条由 \(k\) 标号的余维阶梯。
\end{theorem}
\begin{proof}
由命题 \ref{prop:det-candidate-set-monotone}，候选集合基数随读入单调不增；因此只要长度 \(n\) 后满足 \(|S_n|\ge k\)，则此前所有时刻也都满足 \(|S_t|\ge k\)。于是 \(\mathcal Y_{m,\ge k}\) 的长度-\(n\) 允许前缀数恰为 \(N_{m,\ge k}(n)\)。

而定理 \ref{thm:Ym-ambiguity-spectrum-law} 给出
\[
N_{m,\ge k}(n)=c_{m,\ge k}\rho_{m,\ge k}^{\,n}+O(\rho_{m,\ge k,2}^{\,n}).
\]
在前缀超度量下，长度 \(n\) 柱集正是半径 \(2^{-n}\) 的球；因此完全沿推论 \ref{cor:Ym-amb-hausdorff-codim} 的 covering argument，可得
\[
\dim_{\mathrm H}(\mathcal Y_{m,\ge k})=\frac{\log\rho_{m,\ge k}}{\log 2}.
\]
再代入
\[
\varepsilon_{m,\ge k}=\log 2-\log\rho_{m,\ge k}
\]
即可得到余维公式。证毕。
\end{proof}

\begin{theorem}[扫描预算--余维对偶定理]\label{thm:conclusion-scan-budget-codim-duality}
\leanverified{paper\_conclusion\_scan\_budget\_codim\_duality}
对每个 \(k\ge 2\) 与容忍阈值 \(\eta\in(0,1)\)，定义最小扫描深度
\[
n_{m,\eta}^{(k)}
:=
\min\left\{\,n\ge 1:\ \frac{N_{m,\ge k}(n)}{2^n}\le \eta\,\right\}.
\]
则当 \(\eta\downarrow 0\) 时，
\[
n_{m,\eta}^{(k)}
=
\frac{1}{\varepsilon_{m,\ge k}}\log\frac1\eta
+\frac{1}{\varepsilon_{m,\ge k}}\log c_{m,\ge k}
+O(1).
\]
等价地，
\[
n_{m,\eta}^{(k)}
=
\frac{1}{(\log 2)\,\mathrm{codim}(\mathcal Y_{m,\ge k})}
\log\frac1\eta
+O(1).
\]
因此，扫描时间的主导系数并非某个外加算法常数，而恰是对应阈值歧义集的 Hausdorff 余维；扫描预算与余维在此精确互反。
\end{theorem}
\begin{proof}
由定理 \ref{thm:Ym-ambiguity-spectrum-law}，
\[
\frac{N_{m,\ge k}(n)}{2^n}
=
c_{m,\ge k}e^{-\varepsilon_{m,\ge k}n}(1+o(1)).
\]
令 \(n=n_{m,\eta}^{(k)}\)，则
\[
\eta\asymp c_{m,\ge k}e^{-\varepsilon_{m,\ge k}n},
\]
取对数便得
\[
n
=
\frac{1}{\varepsilon_{m,\ge k}}\log\frac1\eta
+\frac{1}{\varepsilon_{m,\ge k}}\log c_{m,\ge k}
+O(1).
\]
再由定理 \ref{thm:conclusion-threshold-ambiguity-hausdorff-codim-ladder}
\[
\varepsilon_{m,\ge k}=(\log 2)\,\mathrm{codim}(\mathcal Y_{m,\ge k}),
\]
即得第二式。证毕。
\end{proof}

\begin{theorem}[同步可寻址性的唯一指数时间尺度]\label{thm:conclusion-statnull-unique-exponential-address-time-scale}
设
\[
Y_m^{\mathrm{amb}}:=\mathcal Y_{m,\ge 2}\neq\varnothing
\]
且对应阈值子图满足定理 \ref{thm:Ym-ambiguity-spectrum-law} 的“唯一最大谱半径强连通分量且 primitive”假设。记
\[
N_m^{\mathrm{unsync}}(n):=N_{m,\ge 2}(n),
\qquad
\varepsilon_m:=\varepsilon_{m,\ge 2},
\qquad
c_m:=c_{m,\ge 2},
\qquad
\rho_{\mathrm{amb},2}(m):=\rho_{m,\ge 2,2}.
\]
则“何时首次得到可引用地址”在语言层与概率层共享同一唯一指数尺度 \(1/\varepsilon_m\)：
\begin{enumerate}
  \item 未同步长度-\(n\) 前缀的相对密度满足
  \[
  \frac{N_m^{\mathrm{unsync}}(n)}{2^n}
  =
  c_m e^{-\varepsilon_m n}
  +
  O\!\left(\left(\frac{\rho_{\mathrm{amb},2}(m)}{2}\right)^n\right).
  \]
  \item 在最大熵测度下，统计型 \(\mathsf{NULL}\) 的尾概率满足
  \[
  \nu_m(T_m>n)\le C_m e^{-\varepsilon_m n}
  \]
  对某个常数 \(C_m>0\) 成立。
  \item 因而，对 \(N\) 个并行扫描对象，若以组合密度定义未同步前缀的期望个数，则其保持有界的临界时间尺度为
  \[
  n_{\mathrm{sync}}(N,m)=\frac{\log N}{\varepsilon_m}+O(1).
  \]
  更精确地，若
  \[
  n=\frac{\log N+\gamma}{\varepsilon_m},
  \]
  则有
  \[
  N\cdot \frac{N_m^{\mathrm{unsync}}(n)}{2^n}
  =
  c_m e^{-\gamma}+o(1).
  \]
\end{enumerate}
因此，\(\varepsilon_m\) 不只是一个谱参数；它正是“地址首次可引用”的唯一指数时间常数。
\end{theorem}
\begin{proof}
把定理 \ref{thm:Ym-ambiguity-spectrum-law} 应用于 \(k=2\) 即得
\[
N_{m,\ge 2}(n)
=
c_{m,\ge 2}\rho_{m,\ge 2}^{\,n}
 + O(\rho_{m,\ge 2,2}^{\,n}).
\]
两边除以 \(2^n\)，并利用
\[
\varepsilon_m=\log 2-\log\rho_{m,\ge 2},
\qquad
\rho_{\mathrm{amb},2}(m)=\rho_{m,\ge 2,2},
\]
便得到第一条。

第二条正是定理 \ref{thm:Ym-sync-tail} 给出的同步尾指数界，其中指数常数同样等于
\[
\varepsilon_m=\log 2-h_{\mathrm{top}}(Y_m^{\mathrm{amb}}).
\]

对第三条，把
\[
n=\frac{\log N+\gamma}{\varepsilon_m}
\]
代入第一条，可得
\[
N\cdot \frac{N_m^{\mathrm{unsync}}(n)}{2^n}
=
N\left(c_m e^{-\varepsilon_m n}+o(e^{-\varepsilon_m n})\right)
=
c_m e^{-\gamma}+o(1).
\]
于是使该组合期望保持常数级的唯一时间尺度恰为 \((\log N)/\varepsilon_m\)；增减一个常数量级只会对应地放大或压低前述极限常数。证毕。
\end{proof}

\begin{theorem}[清晰度的不可单标量化定理]\label{thm:conclusion-clarity-nonscalarization}
\leanverified{paper\_conclusion\_clarity\_nonscalarization}
若阈值歧义谱并未塌缩，即存在
\[
2\le k<\ell\le |V_m|
\quad\text{使}\quad
\varepsilon_{m,\ge k}<\varepsilon_{m,\ge \ell},
\]
则不存在任何非零标量序列 \(C_m(n)\)，能使所有单调歧义观测都满足
\[
\mathcal A_{\Phi,m}(n)\asymp C_m(n)\qquad (n\to\infty).
\]
更强地，对指标函数
\[
\Phi_k(s):=\mathbf 1_{\{s\ge k\}},
\qquad
\Phi_\ell(s):=\mathbf 1_{\{s\ge \ell\}},
\]
有
\[
\frac{\mathcal A_{\Phi_k,m}(n)}{\mathcal A_{\Phi_\ell,m}(n)}
=
\exp\!\bigl((\varepsilon_{m,\ge \ell}-\varepsilon_{m,\ge k})n+o(n)\bigr)
\longrightarrow +\infty.
\]
因此，一旦残余歧义等级谱真正非平凡，就不存在任何“单一清晰度指标”可以同时代表全部歧义观测；不同层级的残余歧义在指数尺度上彼此不可压缩。
\end{theorem}
\begin{proof}
对上述指标函数，定义直接给出
\[
\mathcal A_{\Phi_k,m}(n)=\frac{N_{m,\ge k}(n)}{2^n},
\qquad
\mathcal A_{\Phi_\ell,m}(n)=\frac{N_{m,\ge \ell}(n)}{2^n}.
\]
再由定理 \ref{thm:Ym-ambiguity-spectrum-law}，
\[
\mathcal A_{\Phi_k,m}(n)
=
c_{m,\ge k}e^{-\varepsilon_{m,\ge k}n}(1+o(1)),
\]
\[
\mathcal A_{\Phi_\ell,m}(n)
=
c_{m,\ge \ell}e^{-\varepsilon_{m,\ge \ell}n}(1+o(1)).
\]
相除即得
\[
\frac{\mathcal A_{\Phi_k,m}(n)}{\mathcal A_{\Phi_\ell,m}(n)}
=
\frac{c_{m,\ge k}}{c_{m,\ge \ell}}
e^{(\varepsilon_{m,\ge \ell}-\varepsilon_{m,\ge k})n}(1+o(1))
\to +\infty.
\]
若存在统一的 \(C_m(n)\) 使两者都与之同阶，则上述比值必须有界，矛盾。证毕。
\end{proof}

\begin{corollary}[“消除统计型 \texorpdfstring{\(\mathsf{NULL}\)}{NULL}”与“压低深层多义度”的非同阶分离]\label{cor:conclusion-statnull-vs-deep-ambiguity-scale-separation}
统计型 \(\mathsf{NULL}\) 与“未同步”等价，即
\[
|S_n(w)|>1
\quad\Longleftrightarrow\quad
w\ \text{在长度 }n\text{ 时仍触发统计型 }\mathsf{NULL}.
\]
若
\[
\varepsilon_{m,\ge 2}<\varepsilon_{m,\ge 3},
\]
并定义
\[
n_{m,\eta}^{(2)}
:=
\min\left\{\,n\ge 1:\ \frac{N_{m,\ge 2}(n)}{2^n}\le \eta\,\right\},
\]
则当 \(\eta\downarrow 0\) 时，
\[
\frac{N_{m,\ge 3}(n_{m,\eta}^{(2)})}{2^{n_{m,\eta}^{(2)}}}
=
\Theta\!\left(\eta^{\varepsilon_{m,\ge 3}/\varepsilon_{m,\ge 2}}\right).
\]
因此，把统计型 \(\mathsf{NULL}\) 压到尺度 \(\eta\) 所需的扫描预算，并不会自动把“三重及以上残余歧义”压到同阶；后者只以幂指数
\[
\varepsilon_{m,\ge 3}/\varepsilon_{m,\ge 2}>1
\]
随 \(\eta\) 衰减。
\end{corollary}
\begin{proof}
由统计型 \(\mathsf{NULL}\) 与未同步的等价及定义，
\[
\frac{N_{m,\ge 2}(n)}{2^n}
\]
正是长度 \(n\) 时仍未同步的最大熵密度。由定理 \ref{thm:conclusion-scan-budget-codim-duality}，
\[
n_{m,\eta}^{(2)}
=
\frac{1}{\varepsilon_{m,\ge 2}}\log\frac1\eta+O(1).
\]
再将其代入
\[
\frac{N_{m,\ge 3}(n)}{2^n}
=
c_{m,\ge 3}e^{-\varepsilon_{m,\ge 3}n}(1+o(1)),
\]
即得
\[
\frac{N_{m,\ge 3}(n_{m,\eta}^{(2)})}{2^{n_{m,\eta}^{(2)}}}
=
\Theta\!\left(\eta^{\varepsilon_{m,\ge 3}/\varepsilon_{m,\ge 2}}\right).
\]
证毕。
\end{proof}

\begin{theorem}[瞬态壳的“周期不可见--清晰度可见”二重分裂定理]\label{thm:conclusion-transient-shell-period-invisible-clarity-visible-splitting}
在一般 sofic 口径下，文稿已证明：歧义壳的瞬态块在确定化邻接矩阵中形成上三角幂零块，因而
\[
\det(I-zA_m^{\mathrm{det}})=\det(I-zA_m^{\mathrm{ess}}),
\]
且周期/Fredholm 数据只感知 singleton skeleton，而不感知 ambiguity shell。与此同时，只要某一级阈值子图 \(H_{m,\ge k}\) 非空并满足定理 \ref{thm:Ym-ambiguity-spectrum-law} 的 Perron--Frobenius 假设，则定理 \ref{thm:conclusion-threshold-ambiguity-hausdorff-codim-ladder} 与定理 \ref{thm:conclusion-scan-budget-codim-duality} 给出
\[
\mathrm{codim}(\mathcal Y_{m,\ge k})
=
\frac{\varepsilon_{m,\ge k}}{\log 2}>0,
\qquad
n_{m,\eta}^{(k)}
=
\frac{1}{\varepsilon_{m,\ge k}}\log\frac1\eta+O(1).
\]
因此，瞬态壳在本体系中出现严格二重分裂：
\begin{enumerate}
\normalfont
  \item 它对周期谱、\(\zeta\)-因子与 determinant 证书不可见；
  \item 它对清晰度预算、残余歧义等级与 Hausdorff 余维完全可见。
\end{enumerate}
换言之，瞬态并非“无谱可见”就等于“可忽略”；它只是从周期通道退场，转而在清晰度通道中留下完整的指数--余维指纹。
\end{theorem}
\begin{proof}
前一半正是定理 \ref{thm:conclusion-ambiguity-shell-zeta-sync-splitting} 的结论：歧义壳作为幂零瞬态块，不进入
\[
\det(I-zA_m^{\mathrm{det}}),\qquad \zeta_{H_m^{\mathrm{det}}}(z),\qquad \{p_n\},\qquad h_{\mathrm{top}}
\]
等周期--熵型不变量。后一半则由定理 \ref{thm:conclusion-threshold-ambiguity-hausdorff-codim-ladder} 与定理 \ref{thm:conclusion-scan-budget-codim-duality} 直接给出：一旦某一级阈值壳非空，其对应残余歧义集便具有严格正的分形余维，并强制出
\[
n_{m,\eta}^{(k)}\sim \varepsilon_{m,\ge k}^{-1}\log(1/\eta)
\]
的对数级扫描预算。二者合并即得所述“周期不可见--清晰度可见”二重分裂。证毕。
\end{proof}

\begin{conjecture}[faithful \texorpdfstring{\(\pi\)}{pi}-formula 的三层证书分裂]\label{conj:conclusion-faithful-pi-formula-three-layer-certificate-splitting}
任一若要在本项目语义中忠实承载 \(\pi\)-通道的公式族，不应只具有单一标量值层，而至少应分裂为三层彼此正交的证书机制
\[
\mathfrak C_{\mathrm{per}}
\oplus
\mathfrak C_{\mathrm{clar}}
\oplus
\mathfrak C_{\det}.
\]
其中：
\begin{enumerate}
\normalfont
  \item \(\mathfrak C_{\mathrm{per}}\) 负责周期轨道、Fredholm 行列式、primitive 计数与其它周期型证书；
  \item \(\mathfrak C_{\mathrm{clar}}\) 负责歧义壳的
  \[
  \varepsilon_{m,\ge k},\qquad \mathrm{codim}(\mathcal Y_{m,\ge k}),\qquad n_{m,\eta}^{(k)}
  \]
  等清晰度量；
  \item \(\mathfrak C_{\det}\) 负责单一行列式势 \(Z_m(t)\) 的端点读出、临界可见证书率 \(r_\varphi\) 与均匀审计律的 AM--GM 熵缺口。
\end{enumerate}
缺失第一层者无法给出周期/Fredholm 完整性；缺失第二层者会把“周期不可见”误判为“语义上可忽略”；缺失第三层者则无法达到一次折叠后的 KL 饱和与端点刚性。定理 \ref{thm:conclusion-transient-shell-period-invisible-clarity-visible-splitting}、定理 \ref{thm:conclusion-single-determinant-potential-three-endpoint-rigidity}、定理 \ref{thm:conclusion-fold-audit-onefold-saturation} 与定理 \ref{thm:conclusion-fold-audit-universal-kl-minimal-idempotent} 提供的正是这一分裂原则的直接证据链。
\end{conjecture}
